\usepackage{amsmath,amsfonts,amsthm,mathtools}
\usepackage{algorithmic}
\usepackage{listings}

\citestyle{acmauthoryear}

\definecolor{mygreen}{rgb}{0,0.4,0}
\definecolor{myblue}{rgb}{0,0,0.7}
\definecolor{myred}{rgb}{0.7,0,0}
\definecolor{mygray}{rgb}{0.5,0.5,0.5}
\definecolor{mymauve}{rgb}{0.58,0,0.82}

\newcommand{\cb}{\color{myblue}}
\usepackage{soul}
% \usepackage{unixode}
\usepackage{wrapfig}

\newcommand*{\mlstinline}[1]{\mbox{\lstinline[mathescape]|#1|}}
\newcommand*{\mlstinlin}[1]{\mbox{\lstinline|#1|}}
\lstset{ 
  backgroundcolor=\color{white},   % choose the background color; you must add \usepackage{color} or \usepackage{xcolor}; should come as last argument
%  basicstyle=\footnotesize,        % the size of the fonts that are used for the code
%  breakatwhitespace=false,         % sets if automatic breaks should only happen at whitespace
%  breaklines=true,                 % sets automatic line breaking
%  captionpos=b,                    % sets the caption-position to bottom
  commentstyle=\color{mygray},    % comment style
%  deletekeywords={...},            % if you want to delete keywords from the given language
%  escapeinside={\%*}{*)},          % if you want to add LaTeX within your code
%  extendedchars=true,              % lets you use non-ASCII characters; for 8-bits encodings only, does not work with UTF-8
%  firstnumber=1000,                % start line enumeration with line 1000
%  frame=single,	                   % adds a frame around the code
%  keepspaces=true,                 % keeps spaces in text, useful for keeping indentation of code (possibly needs columns=flexible)
  keywordstyle=\color{black},       % keyword style
  identifierstyle=\color{myblue},
  language=haskell,                 % the language of the code
  basicstyle=\small\ttfamily,
  morekeywords=[2]{sample,score,Meas},            % if you want to add more keywords to the set
  keywordstyle=[2]\color{myred},       % keyword style
  morekeywords=[3]{Prob,uniform,bernoulli,normal,beta,exponential},            % if you want to add more keywords to the set
  keywordstyle=[3]\color{mygreen},       % keyword style
%  numbers=left,                    % where to put the line-numbers; possible values are (none, left, right)
 % numbersep=5pt,                   % how far the line-numbers are from the code
%  numberstyle=\tiny\color{mygray}, % the style that is used for the line-numbers
 % rulecolor=\color{black},         % if not set, the frame-color may be changed on line-breaks within not-black text (e.g. comments (green here))
  showspaces=false,                % show spaces everywhere adding particular underscores; it overrides 'showstringspaces'
  showstringspaces=false,          % underline spaces within strings only
  showtabs=false,                  % show tabs within strings adding particular underscores
%  stepnumber=2,                    % the step between two line-numbers. If it's 1, each line will be numbered
  stringstyle=\color{mymauve},     % string literal style
 % tabsize=2,	                   % sets default tabsize to 2 spaces
 % title=\lstname                   % show the filename of files included with \lstinputlisting; also try caption instead of title
 escapeinside={(*@}{@*)},
 literate={->}{$\rightarrow\;$}{2}
          {<-}{$\leftarrow\;$}{2}
}

\usepackage[export]{adjustbox}
\usepackage{textcomp}
\usepackage{xcolor}
\usepackage{proof}
\usepackage{stmaryrd}
\usepackage{todonotes}
\usepackage{tikz}
\usepackage{tikz-cd}
\usepackage{cleveref}
\usepackage[caption=false]{subfig}
\DeclareMathOperator*{\colim}{colim}

\newtheorem{theorem}{Theorem}
\newaliascnt{proposition}{theorem}
\newtheorem{proposition}[proposition]{Proposition}
\aliascntresetthe{proposition}
\newaliascnt{definition}{theorem}
\newtheorem{definition}[definition]{Definition}
\aliascntresetthe{definition}
\newaliascnt{lemma}{theorem}
\newtheorem{lemma}[lemma]{Lemma}
\aliascntresetthe{lemma}
\newaliascnt{example}{theorem}
\newtheorem{example}[example]{Example}
\aliascntresetthe{example}
\newaliascnt{corollary}{theorem}
\newtheorem{corollary}[corollary]{Corollary}
\aliascntresetthe{corollary}
%\usepackage{natbib}
\usepackage{microtype}
\DeclareMathOperator{\Aut}{Aut}
\usepackage{mathrsfs}
\usepackage{newunicodechar}
\newunicodechar{：}{\colon}
\newunicodechar{∀}{\forall}
\newunicodechar{∃}{\exists}
\newunicodechar{α}{\alpha}
\newunicodechar{β}{\beta}
\newunicodechar{γ}{\gamma}
\newunicodechar{δ}{\delta}
\newunicodechar{ε}{\epsilon}
\newunicodechar{ζ}{\zeta}
\newunicodechar{η}{\eta}
\newunicodechar{θ}{\theta}
\newunicodechar{ι}{\iota}
\newunicodechar{κ}{\kappa}
\newunicodechar{λ}{\lambda}
\newunicodechar{μ}{\mu}
\newunicodechar{ν}{\nu}
\newunicodechar{ξ}{\xi}
\newunicodechar{ρ}{\rho}
\newunicodechar{σ}{\sigma}
\newunicodechar{τ}{\tau}
\newunicodechar{υ}{\upsilon}
\newunicodechar{φ}{\phi}
\newunicodechar{χ}{\chi}
\newunicodechar{ψ}{\psi}
\newunicodechar{ω}{\omega}
\newunicodechar{Σ}{\Sigma}
\newunicodechar{π}{\pi}
\newunicodechar{Π}{\Pi}
\newunicodechar{Γ}{\Gamma}
\newunicodechar{Δ}{\Delta}
\newunicodechar{⊢}{\vdash}
\newunicodechar{∈}{\in}
\newunicodechar{∉}{\notin}
\newunicodechar{∋}{\ni}
\newunicodechar{∌}{\not\ni}
\newunicodechar{∅}{\emptyset}
\newunicodechar{∪}{\cup}
\newunicodechar{∩}{\cap}
\newunicodechar{∧}{\wedge}
\newunicodechar{∨}{\vee}
\newunicodechar{⊥}{\bot}
\newunicodechar{⊤}{\top}
\newunicodechar{≤}{\leq}
\newunicodechar{≥}{\geq}
\newunicodechar{≠}{\neq}
\newunicodechar{≡}{\equiv}
\newunicodechar{⊆}{\subseteq}
\newunicodechar{⊇}{\supseteq}
\newunicodechar{⊂}{\subset}
\newunicodechar{⊃}{\supset}
\newunicodechar{⟶}{\rightarrow}
\newunicodechar{→}{\to}
\newunicodechar{↪}{\hookrightarrow}
\newunicodechar{⟵}{\leftarrow}
\newunicodechar{↦}{\mapsto}
\newunicodechar{⟹}{\Rightarrow}
\newunicodechar{⟺}{\Leftrightarrow}
\newunicodechar{⊗}{\otimes}
\newunicodechar{×}{\times}
\newunicodechar{∞}{\infty}
\newunicodechar{…}{\ldots}
\newunicodechar{⋯}{\cdots}
\newunicodechar{⋮}{\vdots}
\newunicodechar{⋱}{\ddots}
\newunicodechar{⋰}{\ddots}
\newunicodechar{∧}{\wedge}
\newunicodechar{∨}{\vee}
\newunicodechar{∩}{\cap}
\newunicodechar{∪}{\cup}
\newunicodechar{¬}{\neg}
\newunicodechar{ℙ}{\mathbb{P}}
\newunicodechar{ℕ}{\mathbb{N}}
\newunicodechar{ℝ}{\mathbb{R}}
\newunicodechar{𝕍}{\mathbb{V}}
\newunicodechar{≝}{\defeq}
\newunicodechar{𝒞}{\mathscr{C}}
\newunicodechar{𝒟}{\mathscr{D}}

\newcommand{\hide}[1]{}

\newcommand{\tunit}{\mathsf{unit}}
\newcommand{\treal}{\mathsf{real}}
\newcommand{\tbool}{\mathsf{bool}}
\newcommand{\tvertex}{\mathsf{vertex}}
\newcommand{\tgraph}{\mathsf{graph}}
\newcommand{\tmake}{\mathsf{make}}
\newcommand{\tread}[1]{#1 \#}
\newcommand{\ttrue}{\mathsf{true}}
\newcommand{\tfalse}{\mathsf{false}}
\newcommand{\ite}[3]{\mathsf{if}\,#1\,\mathsf{then}\,#2\,\mathsf{else}\,#3}
\newcommand{\inj}{\mathsf{in}}
\newcommand{\case}[5]{\mathsf{case}\,#1\, \mathsf{of}\,\{\inj_1(#2)\Rightarrow #3;\inj_2(#4)\Rightarrow #5\}}
\newcommand{\casezero}[1]{\mathsf{case}\,#1\,\mathsf{of}\,\{\}}
\newcommand{\casen}[5]{\mathsf{case}\,#1\, \mathsf{of}\,\{#2\Rightarrow #3;\dots;#4\Rightarrow #5\}}
\newcommand{\tnormal}{\mathsf{normal}}
\newcommand{\tbernoulli}{\mathsf{bernoulli}}
\newcommand{\tnew}{\mathsf{new}}
\newcommand{\tedge}{\mathsf{edge}}
\newcommand{\letin}[2]{\mathsf{let}\,#1\,=\,#2\,\mathsf{in}\,}
\newcommand{\letinsqueeze}[2]{\mathsf{let}\,#1{=}#2\,\mathsf{in}\,}
\newcommand{\s}{\,|\,}
\newcommand{\sem}[1]{\llbracket #1\rrbracket}
\newcommand{\semD}[1]{\llbracket #1\rrbracket_{\mathcal{D}}}
\newcommand{\equivD}{\equiv_{\mathcal{D}}}
\newcommand{\RR}{\mathbb{R}}
\newcommand{\NN}{\mathbb{N}}
\newcommand{\VV}{\mathbb{V}}
\newcommand{\AutRado}{\Aut(\mathrm{Rado})}
\newcommand{\id}{\mathrm{id}}
\newcommand{\RadoNom}{\mathbf{RadoNom}}
\newcommand{\Prob}{P}
\newcommand{\inv}{^{\raisebox{0.03ex}{{\textrm{-}}}1}}
\newcommand{\dd}{\mathrm{d}}
\newcommand{\bind}{\mathrel{\scalebox{0.5}[1]{$>\!>=$}}}

\newcommand{\ctxA}{\mathcal{C}}
\newcommand{\ctxE}{\mathcal{E}}
                    \newcommand{\conf}[2]{\langle #1,#2\rangle}

\newcommand{\defeq}{\stackrel {\mathrm{def}}=}
\newcommand{\bernoullisemantics}[2]{\mathrm{Pr}(#1\rightarrow^*#2)}

\newcommand{\Set}{\mathbf{Set}}
\newcommand{\FinStoch}{\mathbf{FinStoch}}
\newcommand{\FinSet}{\mathbf{FinSet}}
\newcommand{\CatA}{\mathbf{C}}
\newcommand{\CatB}{\mathbf{D}}
\newcommand{\CatX}{\mathbf{X}}
\newcommand{\CatC}{\mathbf{A}}
\newcommand{\Meas}{\mathbf{Meas}}
\newcommand{\DistM}{\mathcal{D}}
\newcommand{\Giry}{\mathcal{G}}
\newcommand{\tensor}{\otimes}
\newcommand{\comp}{\circ}
\newcommand{\productmeas}{\otimes}
\newcommand{\marginal}{m}
\newcommand{\fv}{\mathsf{fv}}
\newcommand{\subm}{\mathsf{subm}}

\newcommand{\ie}{\textit{i.e. }}
\newcommand{\eg}{\textit{e.g. }}
\newcommand\restrict[1]{\raisebox{-.5ex}{$|$}_{#1}}
\newcommand{\translation}[1]{\llparenthesis #1\rrparenthesis}
\newcommand{\ErdosRenyi}{Erd\H{o}s--R\'enyi}

\newcommand{\CatAp}{\CatA/\!_\bbase}
\newcommand{\Fam}{\mathsf{Fam}}
\newcommand{\bbase}{\Psi}

